% GENERATED FILE. Edit canonical lessons, then run npm run generate:books.
% canonical-source-hash: fnv1a64:dde96ae11af4e79e
% canonical-lessons: SA-C477-bhramah, SA-C477-vamati, SA-C477-glayati, SA-C477-dahah, SA-C477-jvarati

\chapter{Feeling Unwell}
\label{ch:sa-feeling-unwell}
\hlchaptermodality{477}

\begin{chapteropening}
\textbf{By the end of this chapter:} \emph{I can say that someone felt dizzy or sick.}

Complete the last lesson of chapter 477: I can say that someone felt dizzy or sick.
\end{chapteropening}

\section[\texorpdfstring{\sk{भ्रमः} (bhramaḥ)}{भ्रमः (bhramaḥ)}]{\sk{भ्रमः} (bhramaḥ) — dizziness}

\label{lesson:SA-C477-bhramah}



\begin{quote}
Before the new one: say the Sanskrit for a figure, then the Sanskrit for nourishes.
\end{quote}

\subsection*{You'll want to know: \sk{भ्रमः}}
\textbf{\sk{भ्रमः}} — \emph{bhramaḥ} — \textquotedblleft{}dizziness\textquotedblright{}.

\begin{quote}
\sk{ज्वरेण} \sk{मम} \sk{भ्रमः} \sk{भवति।}
\end{quote}

\begin{quote}
\emph{jvareṇa mama bhramaḥ bhavati.}
\end{quote}

\begin{quote}
The fever makes me dizzy.
\end{quote}

A masculine noun.

\begin{cousinweb}
\textbf{\sk{भ्रमः}} (\emph{bhramaḥ}) comes from the root √bhram, \textquotedblleft{}to wander, to whirl\textquotedblright{}, the root of \textbf{\sk{भ्रमति}}; MW gives \textquotedblleft{}wandering\textquotedblright{} and \textquotedblleft{}giddiness, dizziness\textquotedblright{} (MW).
\end{cousinweb}

\begin{grammarlens}[title={What you've built}]
One more word for describing what you lived and felt.
\end{grammarlens}

\subsection*{Your turn}
\begin{itemize}
\raggedright
  \item \emph{Say it:} \emph{bhramaḥ}
  \item \emph{Say it:} \emph{bhramaḥ}, once more
  \item \emph{Say it:} \emph{jvareṇa mama bhramaḥ bhavati.} — The fever makes me dizzy.
  \item \emph{Recall:} say the Sanskrit for a figure, then the Sanskrit for nourishes, then say \emph{bhramaḥ} again
\end{itemize}

\begin{tcolorbox}[breakable,colback=teal!4,colframe=teal!35!black,title={Before you move on}]
What does \sk{भ्रमः} mean? (\textquotedblleft{}Dizziness\textquotedblright{}.) Say it once more.
\end{tcolorbox}
\section[\texorpdfstring{\sk{वमति} (vamati)}{वमति (vamati)}]{\sk{वमति} (vamati) — vomits}

\label{lesson:SA-C477-vamati}



\begin{quote}
Before the new one: say the Sanskrit for nourishes, then the Sanskrit for dizziness.
\end{quote}

\subsection*{You'll want to know: \sk{वमति}}
\textbf{\sk{वमति}} — \emph{vamati} — \textquotedblleft{}vomits\textquotedblright{}.

\begin{quote}
\sk{बालः} \sk{मलिनं} \sk{जलं} \sk{पिबति}, \sk{अनन्तरं} \sk{वमति।}
\end{quote}

\begin{quote}
\emph{bālaḥ malinaṃ jalaṃ pibati, anantaraṃ vamati.}
\end{quote}

\begin{quote}
The child drinks dirty water and afterwards vomits.
\end{quote}

Past (he …): \sk{सः} \sk{वान्तवान्}, \emph{saḥ vāntavān}.

\begin{cousinweb}
\textbf{\sk{वमति}} (\emph{vamati}) comes from the root √vam, \textquotedblleft{}to vomit\textquotedblright{}, the root of \textbf{\sk{वमनम्}}; MW compares Greek \emph{emeō} and Latin \emph{vomere} (MW).
\end{cousinweb}

\begin{grammarlens}[title={What you've built}]
One more word for describing what you lived and felt.
\end{grammarlens}

\subsection*{Your turn}
\begin{itemize}
\raggedright
  \item \emph{Say it:} \emph{vamati}
  \item \emph{Say it:} \emph{vamati}, once more
  \item \emph{Say it:} \emph{bālaḥ malinaṃ jalaṃ pibati, anantaraṃ vamati.} — The child drinks dirty water and afterwards vomits.
  \item \emph{Recall:} say the Sanskrit for nourishes, then the Sanskrit for dizziness, then say \emph{vamati} again
\end{itemize}

\begin{tcolorbox}[breakable,colback=teal!4,colframe=teal!35!black,title={Before you move on}]
What does \sk{वमति} mean? (\textquotedblleft{}Vomits\textquotedblright{}.) Say it once more.
\end{tcolorbox}
\section[\texorpdfstring{\sk{ग्लायति} (glāyati)}{ग्लायति (glāyati)}]{\sk{ग्लायति} (glāyati) — grows weary}

\label{lesson:SA-C477-glayati}



\begin{quote}
Before the new one: say the Sanskrit for dizziness, then the Sanskrit for vomits.
\end{quote}

\subsection*{You'll want to know: \sk{ग्लायति}}
\textbf{\sk{ग्लायति}} — \emph{glāyati} — \textquotedblleft{}grows weary\textquotedblright{}.

\begin{quote}
\sk{उष्णेन} \sk{वृद्धः} \sk{ग्लायति।}
\end{quote}

\begin{quote}
\emph{uṣṇena vṛddhaḥ glāyati.}
\end{quote}

\begin{quote}
The old man grows faint in the heat.
\end{quote}

Past (he …): \sk{सः} \sk{ग्लानवान्}, \emph{saḥ glānavān}.

\begin{cousinweb}
\textbf{\sk{ग्लायति}} (\emph{glāyati}) comes from the root √glai, \textquotedblleft{}to be wearied, to feel faint\textquotedblright{}; MW's ppp is \emph{glāna} (MW).
\end{cousinweb}

\begin{grammarlens}[title={What you've built}]
One more word for describing what you lived and felt.
\end{grammarlens}

\subsection*{Your turn}
\begin{itemize}
\raggedright
  \item \emph{Say it:} \emph{glāyati}
  \item \emph{Say it:} \emph{glāyati}, once more
  \item \emph{Say it:} \emph{uṣṇena vṛddhaḥ glāyati.} — The old man grows faint in the heat.
  \item \emph{Recall:} say the Sanskrit for dizziness, then the Sanskrit for vomits, then say \emph{glāyati} again
\end{itemize}

\begin{tcolorbox}[breakable,colback=teal!4,colframe=teal!35!black,title={Before you move on}]
What does \sk{ग्लायति} mean? (\textquotedblleft{}Grows weary\textquotedblright{}.) Say it once more.
\end{tcolorbox}
\section[\texorpdfstring{\sk{दाहः} (dāhaḥ)}{दाहः (dāhaḥ)}]{\sk{दाहः} (dāhaḥ) — a burning feeling; inner heat}

\label{lesson:SA-C477-dahah}



\begin{quote}
Before the new one: say the Sanskrit for vomits, then the Sanskrit for grows weary.
\end{quote}

\subsection*{You'll want to know: \sk{दाहः}}
\textbf{\sk{दाहः}} — \emph{dāhaḥ} — \textquotedblleft{}a burning feeling; inner heat\textquotedblright{}.

\begin{quote}
\sk{ज्वरेण} \sk{मम} \sk{देहे} \sk{दाहः} \sk{अस्ति।}
\end{quote}

\begin{quote}
\emph{jvareṇa mama dehe dāhaḥ asti.}
\end{quote}

\begin{quote}
With the fever my body burns.
\end{quote}

A masculine noun.

\begin{cousinweb}
\textbf{\sk{दाहः}} (\emph{dāhaḥ}) comes from the root √dah, \textquotedblleft{}to burn\textquotedblright{}, the root of \textbf{\sk{दहति}}; MW gives \textquotedblleft{}burning, heat\textquotedblright{} and \textquotedblleft{}internal heat, fever\textquotedblright{} (MW).
\end{cousinweb}

\begin{grammarlens}[title={What you've built}]
One more word for describing what you lived and felt.
\end{grammarlens}

\subsection*{Your turn}
\begin{itemize}
\raggedright
  \item \emph{Say it:} \emph{dāhaḥ}
  \item \emph{Say it:} \emph{dāhaḥ}, once more
  \item \emph{Say it:} \emph{jvareṇa mama dehe dāhaḥ asti.} — With the fever my body burns.
  \item \emph{Recall:} say the Sanskrit for vomits, then the Sanskrit for grows weary, then say \emph{dāhaḥ} again
\end{itemize}

\begin{tcolorbox}[breakable,colback=teal!4,colframe=teal!35!black,title={Before you move on}]
What does \sk{दाहः} mean? (\textquotedblleft{}A burning feeling; inner heat\textquotedblright{}.) Say it once more.
\end{tcolorbox}
\section[\texorpdfstring{\sk{ज्वरति} (jvarati)}{ज्वरति (jvarati)}]{\sk{ज्वरति} (jvarati) — runs a fever}

\label{lesson:SA-C477-jvarati}



\begin{quote}
Before the new one: say the Sanskrit for grows weary, then the Sanskrit for a burning feeling.
\end{quote}

\subsection*{You'll want to know: \sk{ज्वरति}}
\textbf{\sk{ज्वरति}} — \emph{jvarati} — \textquotedblleft{}runs a fever\textquotedblright{}.

\begin{quote}
\sk{शीतेन} \sk{बालः} \sk{रात्रौ} \sk{ज्वरति।}
\end{quote}

\begin{quote}
\emph{śītena bālaḥ rātrau jvarati.}
\end{quote}

\begin{quote}
After the cold the child runs a fever at night.
\end{quote}

Past (he …): \sk{सः} \sk{रात्रौ} \sk{ज्वरितवान्}, \emph{saḥ rātrau jvaritavān}.

\begin{cousinweb}
\textbf{\sk{ज्वरति}} (\emph{jvarati}) comes from the root √jvar, \textquotedblleft{}to be hot, to be feverish\textquotedblright{}, the root of \textbf{\sk{ज्वरः}} (MW).
\end{cousinweb}

\begin{grammarlens}[title={What you've built}]
One more word for describing what you lived and felt.
\end{grammarlens}

\subsection*{Your turn}
\begin{itemize}
\raggedright
  \item \emph{Say it:} \emph{jvarati}
  \item \emph{Say it:} \emph{jvarati}, once more
  \item \emph{Say it:} \emph{śītena bālaḥ rātrau jvarati.} — After the cold the child runs a fever at night.
  \item \emph{Recall:} say the Sanskrit for grows weary, then the Sanskrit for a burning feeling, then say \emph{jvarati} again
\end{itemize}

\begin{tcolorbox}[breakable,colback=teal!4,colframe=teal!35!black,title={Before you move on}]
What does \sk{ज्वरति} mean? (\textquotedblleft{}Runs a fever\textquotedblright{}.) Say it once more.
\end{tcolorbox}
